\documentclass[11pt]{article}
\usepackage[margin=1in]{geometry}
\usepackage[T1]{fontenc}
\usepackage{lmodern}
\usepackage{amsmath,amssymb,amsthm,mathtools}
\usepackage{enumitem}
\usepackage{hyperref}

\newtheorem{theorem}{Theorem}
\newtheorem{lemma}{Lemma}
\newtheorem{remark}{Remark}
\newcommand{\E}{\mathcal E_p}
\newcommand{\R}{\mathbb R}
\newcommand{\Sone}{\mathbb S^1}
\newcommand{\ip}[2]{\left\langle #1,#2\right\rangle}

\title{Subquadratic $p$-Elastic Flow Near Nonzero Curvature:\\
An $H^4\!\to L^2$ Convergence Extension}
\author{Run fa\_banach\_001, agent\_lane\_14}
\date{11 August 2026}

\begin{document}
\maketitle

\section*{Status and source question}

\textbf{Status:} \texttt{candidate\_partial\_likely\_valid}.

Pozzetta's \emph{Convergence of elastic flows of curves into manifolds},
arXiv:2007.00582 \cite{Pozzetta}, studies
\[
  \E(\gamma)=\int_{\Sone}\left(1+\frac1p|k|^p\right)\,ds
\]
and its geometric fourth-order flow for $p\ge2$.  On PDF page~5 the paper
states that understanding the behavior for $p\in[1,2)$ remains open.  This
packet proves a conditional convergence theorem for the full subquadratic
range $1<p<2$ whenever the curvature stays away from zero.  It also explains
why this does not yet settle the broad source problem.

Let $p'=p/(p-1)$ and write the normal gradient as
\begin{equation}\label{eq:gradient}
 G_p(\gamma)=\nabla_s^2\bigl(|k|^{p-2}k\bigr)
 +\frac1{p'}|k|^p k-k
 +R\bigl(|k|^{p-2}k,\tau\bigr)\tau .
\end{equation}
The flow is $\partial_t\gamma=-G_p(\gamma)$, up to tangential
reparametrization.

\section*{New result}

\begin{theorem}[subquadratic nonzero-curvature regime]\label{thm:main}
Let $1<p<2$, and let $(M,g)$ be a complete real-analytic Riemannian manifold
with real-analytic metric.
\begin{enumerate}[label=\textup{(\roman*)}]
\item If a smooth immersed closed curve $\gamma_0$ satisfies
$\min_{\Sone}|k_{\gamma_0}|>0$, then \eqref{eq:gradient} has a unique smooth
short-time solution, modulo reparametrization, and its curvature remains
nonzero for a short time.
\item If $\gamma_*$ is a smooth critical point of $\E$ with
$\min_{\Sone}|k_{\gamma_*}|>0$, there are $C,\rho>0$ and
$\theta\in(0,1/2]$ such that every $H^4$ normal graph $\gamma_v$ over
$\gamma_*$ with $\|v\|_{H^4}<\rho$ satisfies
\begin{equation}\label{eq:LS}
 |\E(\gamma_v)-\E(\gamma_*)|^{1-\theta}
 \le C\,\|G_p(\gamma_v)\|_{L^2(ds_{\gamma_v})}.
\end{equation}
\item Let $\gamma:[0,\infty)\times\Sone\to M$ be a global smooth solution
with $|k_\gamma|>0$.  Assume that there are times $t_j\to\infty$, ambient
isometries $I_j$, reparametrizations $\phi_j$, and a smooth critical curve
$\gamma_*$ with $\min|k_{\gamma_*}|>0$ such that
\[
 I_j\circ\gamma(t_j,\phi_j(\cdot))\longrightarrow\gamma_*
 \quad\text{in }C^m(\Sone)\quad\text{for every }m.
\]
Then the flow converges smoothly to a critical curve as $t\to\infty$, up to
reparametrization.  Thus the source theorem's ``sub-convergence implies full
convergence'' conclusion extends to $1<p<2$ in the nonzero-curvature regime.
\end{enumerate}
\end{theorem}

\section*{The exponent obstruction and the Hilbert repair}

The source realizes the first variation in $L^{p'}$ against $L^p$ and, in
the convergence argument, uses
\[
 \|V\|_{L^2}\ge c\|V\|_{L^{p'}}.
\]
This is available on a finite-length curve when $p'\le2$, equivalently
$p\ge2$.  For $1<p<2$ one has $p'>2$ and the inclusion reverses.  Moreover,
an $H^{4,p}$ variation naturally places the fourth-order expression in
$L^p$, not in $L^{p'}$.  Merely changing the stated range of $p$ in the
source proof therefore does not work.

The repair is to use normal fields in the Hilbert pair
\[
 V=H^4(\Sone;N\gamma_*),\qquad Z=L^2(\Sone;N\gamma_*),
\]
and represent the first variation by its fixed-density $L^2$ gradient.  This
produces \eqref{eq:LS} directly in precisely the norm occurring in energy
dissipation.

\section*{Proof of the local analytic and elliptic statements}

Put $e=k/|k|$ and define, on the normal bundle,
\begin{equation}\label{eq:A}
 A_k=|k|^{p-2}\bigl(I+(p-2)e\otimes e\bigr).
\end{equation}
For a normal vector $\xi=\xi_\perp+\xi_\parallel$, with $\xi_\parallel$
parallel to $k$, direct calculation gives
\begin{equation}\label{eq:eigs}
 \ip{A_k\xi}{\xi}
 =|k|^{p-2}\left(|\xi_\perp|^2+(p-1)|\xi_\parallel|^2\right).
\end{equation}
Thus $A_k$ is positive definite for every $p>1$.  On an annulus
$0<c\le |k|\le C_0$, its ellipticity constants are uniform.  Notice that
the possibly negative coefficient $p-2$ in \eqref{eq:A} causes no problem;
the longitudinal eigenvalue is $(p-1)|k|^{p-2}>0$.

Choose a normal-graph chart $v\mapsto\gamma_v=\exp_{\gamma_*}v$.  Since
$H^4(\Sone)$ embeds into $C^3$, small $H^4$ graphs remain immersed and obey
$|k_{\gamma_v}|\ge c/2$.  Curvature is an analytic function of the graph and
its first two derivatives.  The scalar function $r\mapsto r^{p/2}$ is real
analytic for $r>0$, so both
\[
 v\longmapsto \E(\gamma_v)\in\R,
 \qquad
 v\longmapsto \mathcal M(v)\in L^2
\]
are analytic; here $\mathcal M(v)$ is \eqref{eq:gradient}, transported to
$N\gamma_*$ and multiplied by the analytic density change.  The second
variation at the critical point has the form
\begin{equation}\label{eq:hessian}
 L v=\nabla_s^2\bigl(A_k\nabla_s^2v\bigr)+Kv,
 \qquad L:H^4\longrightarrow L^2,
\end{equation}
where $K$ contains at most three derivatives of $v$.  Hence $K:H^4\to L^2$
is compact by Rellich's theorem.

Let $P=\nabla_s^2(A_k\nabla_s^2)$.  Equations \eqref{eq:eigs} and standard
one-dimensional interpolation give the coercive estimate
\[
 \|v\|_{H^2}^2
 \le C\left(\int_{\Sone}\ip{A_k\nabla_s^2v}{\nabla_s^2v}\,ds
 +\|v\|_{L^2}^2\right).
\]
Lax--Milgram followed by fourth-order elliptic regularity shows that
$P+I:H^4\to L^2$ is an isomorphism.  Therefore $P$, and then $L=P+K$, is
Fredholm of index zero.  Chill's analytic gradient theorem
\cite[Corollary 3.11]{Chill}, in the $H^4\to L^2$ realization, now gives
\eqref{eq:LS}.  This proves part~\textup{(ii)}.

For part~\textup{(i)}, fix a standard tangential gauge.  In a local chart
the fourth-order principal symbol is, up to a positive power of
$|\partial_x\gamma|$, $A_k\xi^4$.  Formula \eqref{eq:eigs} makes the gauged
system uniformly parabolic at $\gamma_0$.  Classical quasilinear
fourth-order parabolic theory gives a unique smooth gauged solution.  The
usual geometric-uniqueness argument removes the gauge, and continuity keeps
$|k|$ positive for a short time.

\section*{Proof of conditional full convergence}

Only the exponent-sensitive part of the source's convergence argument needs
replacement.  We give the trapping mechanism explicitly.

After choosing a sufficiently late subsequence time, applying its isometry,
and using a constant-speed parametrization, start the flow at a curve
$\varepsilon$-close to $\gamma_*$ in every fixed $C^m$ norm.  Let $T$ be the
maximal time for which the reparametrized flow remains in the normal-graph
$H^4$ ball where \eqref{eq:LS} holds and $|k|\ge\frac12\min|k_{\gamma_*}|$.
On $[0,T)$ set
\[
 H(t)=\bigl(\E(\gamma(t))-\E(\gamma_*)\bigr)^\theta.
\]
The normal flow satisfies the exact dissipation identity
\begin{equation}\label{eq:diss}
 -\frac{d}{dt}\E(\gamma(t))=\|G_p(\gamma(t))\|_{L^2(ds)}^2.
\end{equation}
Combining \eqref{eq:LS} and \eqref{eq:diss} yields
\begin{equation}\label{eq:length}
 -H'(t)
 =\theta\frac{\|G_p(\gamma(t))\|_2^2}
 {\bigl(\E(\gamma(t))-\E(\gamma_*)\bigr)^{1-\theta}}
 \ge c\|G_p(\gamma(t))\|_2.
\end{equation}
Consequently the $L^2$ length of the normal trajectory before exit is at
most $C H(0)=O(\varepsilon^\theta)$.

For completeness, the parabolic estimate used to turn this path-length
bound into trapping also survives for $1<p<2$.  After differentiating the
curvature equation $m$ times, its leading integrated term is
\[
 -\int_{\Sone}
 \ip{A_k\nabla_s^{m+2}k}{\nabla_s^{m+2}k}\,ds.
\]
By \eqref{eq:eigs} this controls the highest derivative with constant
proportional to $\min\{1,p-1\}$.  All remaining coefficients are smooth and
bounded on the curvature annulus.  One-dimensional interpolation absorbs
the intermediate derivatives, giving the same uniform high-Sobolev bounds
while the curve remains in the chosen $H^4$ neighborhood as in the source's
parabolic-estimate lemma.

In constant-speed gauge, the change of parametrization is controlled by
$\int|k||G_p|\,ds$.  The uniform curvature bound and \eqref{eq:length}
therefore give
\[
 \|\widehat\gamma(t)-\gamma_*\|_{L^2}
 \le C\varepsilon^\theta\qquad(0\le t<T).
\]
Interpolating this small $L^2$ distance against the uniform high-Sobolev
bounds yields
\[
 \|\widehat\gamma(t)-\gamma_*\|_{H^4}
 \le C\varepsilon^\alpha
\]
for some $\alpha>0$ independent of $T$.  Taking $\varepsilon$ small makes
the right side strictly smaller than the exit radius, contradicting finite
maximality.  Hence $T=\infty$.

Equation \eqref{eq:length} now gives finite total $L^2$ path length.  The
uniform parabolic bounds and interpolation upgrade the resulting $L^2$
Cauchy property to convergence in every $C^m$ norm.  The limit is critical
by \eqref{eq:diss} and continuity of $G_p$.  Finally, the source's compactness
argument for the sequence of ambient isometries is independent of the value
of $p$ once trapping and finite $L^2$ path length are known; it identifies a
single limiting critical curve for the original flow.  This proves
part~\textup{(iii)}.

\section*{Upgrade audit and remaining obstruction}

The result was pushed along four further directions.
\begin{itemize}
\item \textbf{Curvature zeros.}  For $1<p<2$, the map
$z\mapsto|z|^{p-2}z$ is not differentiable at $z=0$, and its linearization
is singular there.  Analytic gradient theory and uniform classical
parabolicity both fail.  A nonsmooth/subgradient theory is required.
\item \textbf{Global existence and sub-convergence.}  The local theorem does
not prevent curvature from reaching zero and supplies no global continuation
estimate.  These are the principal missing inputs for a complete solution of
the source's broad problem.
\item \textbf{The endpoint $p=1$.}  The longitudinal eigenvalue in
\eqref{eq:eigs} becomes zero, and the curvature integrand is nonsmooth.  The
present proof has a sharp structural endpoint barrier.
\item \textbf{Quantitative confinement and Huisken's half-plane conjecture.}
The path-length estimate begins only after entry into a local
\L ojasiewicz--Simon neighborhood and depends on the limiting critical curve.
It gives neither an initial-data-only confinement radius nor a half-plane
principle.
\end{itemize}

Bounded title, author, exact-phrase, and arXiv searches through 11 August
2026 found no prior fourth-order subquadratic nonzero-curvature convergence
theorem.  The related arXiv:1811.06608 concerns a different second-order
inextensible flow, while arXiv:2104.03570 treats a modified weak flow for
$p\ge2$.  A 2026 survey/thesis still lists the broad subquadratic and
half-plane questions as open.  Novelty confidence is moderate because this
packet extends the source machinery rather than introducing a new global
estimate.

\begin{thebibliography}{9}
\bibitem{Pozzetta}
M. Pozzetta,
\emph{Convergence of elastic flows of curves into manifolds},
arXiv:2007.00582 (2020; revised 2021).

\bibitem{Chill}
R. Chill,
\emph{On the \L ojasiewicz--Simon gradient inequality},
J. Funct. Anal. \textbf{201} (2003), 572--601.
\end{thebibliography}

\end{document}
